\documentclass[10pt,a4paper]{amsart}
\usepackage[margin=19mm]{geometry}
\usepackage{amsmath,amssymb,mathtools}
\usepackage[scr=rsfs,cal=euler]{mathalfa}
\usepackage{xfrac}
\usepackage[unicode,hidelinks,bookmarks=false]{hyperref}
\newcommand{\ie}{i.e.\ }
\newcommand{\cf}{cf.\ }
\newcommand{\resp}{respectively\ }
\newcommand{\Subsection}{\subsection}
\providecommand{\nobreakdash}{\nobreak\hbox{-}}
\setlength{\emergencystretch}{2em}
\multlinegap0pt
\usepackage[shortlabels]{enumitem}
\usepackage{mathtools}
\usepackage{amssymb}
\usepackage{bm}
\usepackage{algorithm}\usepackage{algpseudocode}
\usepackage{xcolor}
\newtheorem{lemma}{Lemma}
\newcommand{\KD} {\texttt{KD}}
\newcommand{\KI} {\texttt{KI}}
\title{Tail Distribution of Regret in Optimistic Reinforcement Learning}
\author{Sajad Khodadadian, Mehrdad Moharrami}
\date{Source: arXiv:2511.18247v3 (CC BY 4.0)}
\begin{document}
\maketitle

\noindent\emph{Source and licence.} Tail Distribution of Regret in Optimistic Reinforcement Learning. Version arXiv:2511.18247v3 (CC BY 4.0), distributed by arXiv under the Creative Commons Attribution 4.0 International licence, http://creativecommons.org/licenses/by/4.0/, as recorded in the arXiv licence metadata for this exact version and shown on the abstract page https://arxiv.org/abs/2511.18247v3. Full licence notice: \texttt{licenses/arxiv-2511.18247-CC-BY-4.0.txt}.\par\smallskip

\noindent\textit{Editorial scope.} Counted unit: the article's lemma lem:sumUpperBound-UCBVI (source0.tex lines 1083--1095). The article's complete original proof of it is delivered with it (source0.tex lines 1096--1147, one \texttt{proof} environment of 52 lines). No mathematics is written, repaired or re-derived: the statement and the proof are the article's own lines, in the article's own notation, and no retained line is substituted or re-flowed.  No further source-local context is needed: every cross-reference of every retained line resolves inside the two delivered spans.  Nothing was omitted from any retained span, so the package carries no omission record.  No retained line contains a citation, so the wrapper carries no bibliography and no bibliography entry.  No retained line contains a figure environment, an image-inclusion command or drawing code, so no image is delivered and the package declares no asset dependency.  Editorial formatting only, in addition: an AMS \texttt{amsart} preamble that restates the article's own declarations and macros used by the retained lines, namely the environments \texttt{lemma} on separate counters, together with the standard packages those lines need.  The retained environments print in this document as Lemma 1 in order of appearance, whereas the article prints the same items with the numbering of its own full document.  The complete native source of the article as distributed in the arXiv e-print for this exact version is bundled unchanged as \texttt{sources/189640/source0.tex}.
\begin{lemma}\label{lem:sumUpperBound-UCBVI}
The following bound holds:
\begin{align*}
    \sum_{n=1}^K 
        &\exp\left(-  2n  \left(\frac{b_h^{n \vee \lceil K^\gamma\rceil}(n)}{H-h-1} \right)^2\right) \leq 2^{-S}\times
    \begin{cases}
        \displaystyle K \exp\left(-2\mu { K^{\alpha}}\right),
        & \text{\KD}, \\[2ex]
        \displaystyle K \exp\left(-2\mu  K^{\gamma\alpha}\right),
        & \text{\KI}.
    \end{cases}
\end{align*}
\end{lemma}

\begin{proof}[Proof of Lemma \ref{lem:sumUpperBound-UCBVI}]
Recall that the function $b^k_h(\cdot)$ is defined as
\[
b^k_h(n)
=
\begin{cases}
(H-h-1)\sqrt{\dfrac{0.5 S\ln 2+\mu K^\alpha}{n}}, & \text{\KD},\\[2ex]
(H-h-1)\sqrt{\dfrac{0.5 S\ln 2+\mu (k+1)^\alpha}{n}}, & \text{\KI},
\end{cases}
\qquad
b^k_h(0)\coloneqq\infty.
\]
Let $S_K$ denote the summation in Lemma~\ref{lem:sumUpperBound-UCBVI}, i.e.,
\begin{align*}
    S_K \coloneqq \sum_{n=1}^{K} \exp\left(-  2n  \left(\frac{b_h^{n \vee \lceil K^\gamma\rceil}(n)}{H-h-1} \right)^2\right).
\end{align*}
The analysis depends on the choice of $b^k_h(\cdot)$.

\begin{itemize}[leftmargin=0pt]
    \item {\bf \KD:} We have
        \begin{align*}
        S_K &\leq \sum_{n=1}^{K} \exp\left(-S\ln2-2\mu  K^{\alpha}\right) = 2^{-S} K \exp\left({- 2\mu K^{\alpha}}\right) 
        \end{align*}
    \item {\bf \KI:} Define $S^{(1)}_K$ and $S^{(2)}_K$ as follows:
    \begin{align*}
        S^{(1)}_K \coloneqq \sum_{n=1}^{\lceil K^\gamma\rceil} \exp\left(- 2n  \left(\frac{b_h^{\lceil K^\gamma\rceil}(n)}{H-h-1} \right)^2\right),
        \qquad
        S^{(2)}_K \coloneqq \sum_{n=\lceil K^\gamma\rceil + 1}^{K} \exp\left(- 2n \left(\frac{b_h^{n}(n)}{H-h-1} \right)^2\right).
    \end{align*}
    Notice that $S_K = S^{(1)}_K + S^{(2)}_K,$ and that
        \begin{align*}
            b_h^{\lceil K^\gamma\rceil}(n) = (H-h-1)\sqrt{\frac{0.5 S\ln2 + \mu(\lceil K^\gamma\rceil+1)^{\alpha}}{n}},
            \quad
            b_h^n(n) = (H-h-1)\sqrt{\frac{0.5 S\ln2 + \mu( n+1 )^{\alpha}}{n}}.
        \end{align*}
        
        We proceed by bounding $S_K^{(1)}$ and $S_K^{(2)}$. We begin with $S_K^{(1)}$, for which we have
        \begin{align*}
        S_K^{(1)}
        &= \sum_{n=1}^{\lceil K^\gamma\rceil}\exp\left(-2n \left(\frac{0.5S\ln2 + \mu(\lceil K^\gamma\rceil+1)^{\alpha}}{n}\right)\right) \leq 2^{-S}\lceil K^\gamma\rceil \exp\left(-2\mu K^{\gamma\alpha}\right) 
        \end{align*}
    \end{itemize}

We next bound $S_K^{(2)}$:
\begin{align*}
S^{(2)}_K
&=\sum_{n=\lceil K^\gamma\rceil + 1}^{K} \exp\left(- 2n \left(\frac{b_h^{n}(n)}{H-h-1} \right)^2\right)\\
&\le\sum_{n= \lceil K^\gamma\rceil+1}^{K}2^{-S}\exp\left(-2\mu  n^{\alpha}\right) \\
&\leq 2^{-S} (K - \lceil K^\gamma\rceil)\exp\left(-2\mu  K^{\gamma\alpha}\right).
\end{align*}
Putting everything together, we obtain the desired result.
\end{proof}

\end{document}
